% ============================================================================
\section{Does Loss Design Change Update Sparsity?}
\label{sec:density}
% ============================================================================

Our third study examines vocabulary-level supervision density: how many
vocabulary alternatives enter the loss at each student-generated prefix.
The endpoints are sampled-token policy gradients (PG) and full-vocabulary
KL; top-$K$ losses provide an intermediate practical choice. This concerns
vocabulary coverage, not how many response positions receive supervision.
We compare these choices in single-teacher and multi-teacher settings
(Table~\ref{tab:density_controls}). Prior work already observes concentrated
OPD parameter changes under dense feedback \citep{yu2026denseupdates}.
We ask whether these loss choices produce materially different
update sparsity under otherwise matched configurations, without assuming
that sampled supervision is sparser or reduces interference.

\begin{table}[t]
\centering
\small
\begin{tabular}{lll}
\toprule
Setting & Sampled-token PG & Vocabulary-level KL \\
\midrule
Single teacher & Single-teacher setting & Full-vocabulary / top-$K$ \\
Multiple teachers & Joint MOPD setting & Full-vocabulary / top-$K$ \\
\bottomrule
\end{tabular}
\caption{Primary loss-design comparison. Both settings include full-vocabulary
local probes; online PG/top-$K$ runs provide cumulative sparsity and
capability comparisons.}
\label{tab:density_controls}
\end{table}

Full-vocabulary loss is evaluated on a small common prefix bank in both
settings, producing a reference gradient and proposed optimizer step.
It is part of the primary density measurement, not only a follow-up
triggered by a PG/top-$K$ difference. Cumulative training comparisons
initially use PG and top-$K$; full-vocabulary online training is subject
to measured feasibility and is not implied by a local probe.

Within each setting, paired runs share the student initialization, teachers,
prompt schedule, length limits, task normalization, optimizer configuration,
and training steps. We fix and document $K$, support selection, and probability
normalization for the chosen top-$K$ implementation
\citep{ma2026mopdmultiteacheronpolicydistillation}. Online rollouts follow each
run's evolving student; generated-token exposure and compute are reported
alongside matched-step comparisons.

We measure actual optimizer-step sparsity and cumulative parameter changes
at identical checkpoints, using the same FP32 measurement procedure and
thresholds across conditions. Update norms and scale-normalized concentration
provide context for threshold-visible sparsity. Repeated seeds yield effect
estimates and confidence intervals. Held-out task scores establish whether
a sparser update accompanies comparable learning; sparsity alone does not
demonstrate better capability integration.

At a fixed prefix, fresh sampled-token PG with detached log-ratio coefficients
has the same expected gradient as full-vocabulary reverse KL, a known identity
\citep{yu2026denseupdates,oh2026vopd}. This does not equate practical top-$K$
objectives, length-weighted original rollout samples, nonlinear Adam updates,
or long-term parameter supports.

If confidence intervals fall within a prespecified materiality margin, we
will report no material sparsity difference within the tested regime and
stop this line of analysis; intervals too wide to resolve that margin are
inconclusive. Only a stable, meaningful gap would motivate an optional
follow-up separating sampling noise from truncation bias, for example using
a variance-reduced sampled estimator on the common reference bank.
These mechanism controls are conditional extensions, not required components
of the primary study, and no preference for sampled PG is assumed.
